% Canonical LaTeX source; imported and revised from the legacy Markdown.
\chapter{Serie numeriche}
\label{chap:17}

\begin{obiettivocapitolo}{scegliere un criterio conclusivo per la convergenza di una serie}
\prioritaalta. Il primo controllo è sempre $a_n\to0$; se fallisce, la serie diverge. Se vale, non conclude nulla da solo.
\end{obiettivocapitolo}
\begin{sceltametodo}{albero decisionale per le serie}
\begin{enumerate}
\item \textbf{Riconosci serie notevoli:} geometrica, telescopica, armonica/generalizzata.
\item \textbf{Termini non negativi:} confronto diretto o asintotico; rapporto/radice per fattoriali ed esponenziali; integrale o condensazione per forme lente.
\item \textbf{Segni alterni:} prova convergenza assoluta; se fallisce, usa Leibniz verificando monotonia e limite a zero.
\item \textbf{Oscillazioni non alternate semplici:} considera Dirichlet o Abel.
\item \textbf{Serie di potenze:} trova il raggio e studia separatamente ogni estremo.
\end{enumerate}
\end{sceltametodo}
\begin{proceduraesame}{applicare un criterio}
\begin{enumerate}
\item Scrivi chiaramente $a_n$ e l'indice iniziale.
\item Verifica le ipotesi del criterio scelto.
\item Calcola il limite o il confronto e traduci il risultato in convergenza/divergenza.
\item Se il criterio è inconcludente, dichiaralo e passa a un altro; non dedurre una conclusione.
\item Se richiesto, distingue convergenza assoluta, condizionata e stima del resto.
\end{enumerate}
\end{proceduraesame}
\begin{errorefrequente}{rapporto e radice}
Quando il limite vale $1$, i criteri del rapporto e della radice sono inconcludenti. Non indicano convergenza né divergenza.
\end{errorefrequente}
\begin{controllofinale}{serie}
Confronta la conclusione con il termine generale e con una serie di riferimento; specifica sempre il comportamento agli estremi nei problemi parametrici.
\end{controllofinale}

\[
\sum a_n\text{ e }\sum b_n\text{ hanno lo stesso carattere}
\Longleftrightarrow
\bigl(\text{entrambe convergenti}\bigr)
\lor
\bigl(\text{entrambe divergenti}\bigr).
\]

\section{Definizioni e proprietà}\label{definizioni-e-proprietuxe0}

Per \(a_n\in\mathbb K\), con \(\mathbb K\in\{\mathbb R,\mathbb C\}\):

\[
\displaystyle
\sum_{n=n_0}^{\infty}a_n,
\qquad
s_N=\sum_{n=n_0}^{N}a_n,
\]

\[
\displaystyle
\sum_{n=n_0}^{\infty}a_n=A
\quad\Longleftrightarrow\quad
s_N\to A.
\]

Se la serie converge a \(A\), il \textbf{resto dopo il termine di indice \(N\)} è

\[
R_N=A-s_N=\sum_{n=N+1}^{\infty}a_n,
\qquad
R_N\to0.
\]

Invarianza rispetto a un numero finito di termini, per \(N_0\ge n_0\):

\[
\displaystyle
\sum_{n=n_0}^{\infty}a_n
\ \text{e}\
\sum_{n=N_0}^{\infty}a_n
\quad
\text{hanno lo stesso carattere},
\qquad N_0\ge n_0.
\]

Linearità:

\[
\displaystyle
\sum a_n=A,\quad \sum b_n=B
\quad\Longrightarrow\quad
\sum(\alpha a_n+\beta b_n)=\alpha A+\beta B.
\]

Per \(a_n\ge0\) definitivamente:

\[
\displaystyle
\sum a_n\ \text{converge}
\quad\Longleftrightarrow\quad
(s_N)\ \text{è limitata superiormente},
\]

\[
\displaystyle
(s_N)\ \text{non limitata}
\quad\Longrightarrow\quad
s_N\to+\infty.
\]

Teoria: \href{https://ingegnerismo.it/matematica/serie-numerica/}{serie numerica}.

\section{Condizione necessaria e criterio di Cauchy}\label{condizione-necessaria-e-criterio-di-cauchy}

\[
\displaystyle
\sum a_n\ \text{converge}
\quad\Longrightarrow\quad
a_n\to0.
\]

\[
\displaystyle
a_n\not\to0
\quad\Longrightarrow\quad
\sum a_n\ \text{diverge}.
\]

Il converso della condizione necessaria è falso:

\[
\frac1n\to0,
\qquad
\sum_{n=1}^{\infty}\frac1n\text{ diverge}.
\]

\[
\displaystyle
\sum a_n\ \text{converge}
\quad\Longleftrightarrow\quad
\forall\varepsilon>0\ \exists N\ \forall p,q:
\ q\ge p\ge N
\Longrightarrow
\left\lvert\sum_{n=p}^{q}a_n\right\rvert<\varepsilon.
\]

Teoria: \href{https://ingegnerismo.it/matematica/condizione-necessaria-convergenza/}{condizione necessaria di convergenza} e \href{https://ingegnerismo.it/matematica/criterio-di-cauchy/}{criterio di Cauchy}.

\section{Serie notevoli}\label{serie-notevoli}

Serie geometrica, \(q\ne1\):

\[
\displaystyle
\sum_{n=0}^{N}q^n
=
\frac{1-q^{N+1}}{1-q}.
\]

Per \(q=1\):

\[
\displaystyle
\sum_{n=0}^{N}1=N+1.
\]

\[
\displaystyle
\sum_{n=0}^{\infty}q^n
=
\begin{cases}
\dfrac1{1-q}, & \lvert q\rvert<1,\\[4pt]
\text{diverge}, & \lvert q\rvert\ge1.
\end{cases}
\]

Per \(\lvert q\rvert<1\):

\[
\displaystyle
R_N=\sum_{n=N+1}^{\infty}q^n
=
\frac{q^{N+1}}{1-q},
\]

\[
|R_N|=\frac{|q|^{N+1}}{|1-q|}.
\]

Per garantire \(|R_N|\le\varepsilon\) si sceglie \(N\) tale che

\[
|q|^{N+1}\le\varepsilon|1-q|.
\]

Serie armonica e \(p\)-serie:

\[
\displaystyle
\sum_{n=1}^{\infty}\frac1n
\quad\text{diverge},
\]

\[
\displaystyle
\sum_{n=1}^{\infty}\frac1{n^p}\text{ converge}
\Longleftrightarrow p>1,
\]

\[
\displaystyle
\sum_{n=1}^{\infty}\frac1{n^p}\text{ diverge}
\Longleftrightarrow p\le1.
\]

Serie logaritmica, per \(p,q\in\mathbb R\); il carattere è dato da:

\[
\displaystyle
\sum_{n=2}^{\infty}\frac1{n^p(\ln n)^q}\text{ converge}
\Longleftrightarrow
p>1\ \lor\ (p=1\ \land\ q>1),
\]

\[
\displaystyle
\sum_{n=2}^{\infty}\frac1{n^p(\ln n)^q}\text{ diverge}
\Longleftrightarrow
p<1\ \lor\ (p=1\ \land\ q\le1).
\]

Serie telescopica:

\[
\displaystyle
\sum_{n=1}^{N}(b_n-b_{n+1})
=
b_1-b_{N+1},
\]

\[
\displaystyle
\sum_{n=1}^{\infty}(b_n-b_{n+1})
\ \text{converge}
\quad\Longleftrightarrow\quad
b_n\to L\in\mathbb K,
\]

\[
\displaystyle
\sum_{n=1}^{\infty}(b_n-b_{n+1})
=
b_1-L.
\]

Serie armonica alternata:

\[
\displaystyle
\sum_{n=1}^{\infty}\frac{(-1)^{n-1}}n=\ln2.
\]

Serie di Mengoli:

\[
\frac1{n(n+1)}=\frac1n-\frac1{n+1},
\]

\[
\sum_{n=1}^{N}\frac1{n(n+1)}=1-\frac1{N+1},
\qquad
\sum_{n=1}^{\infty}\frac1{n(n+1)}=1.
\]

Per \(|q|<1\):

\[
\sum_{n=1}^{\infty}nq^{n-1}=\frac1{(1-q)^2},
\qquad
\sum_{n=1}^{\infty}nq^n=\frac{q}{(1-q)^2}.
\]

Teoria e casi: \href{https://ingegnerismo.it/matematica/serie-notevoli/}{serie notevoli} e \href{https://ingegnerismo.it/matematica/serie-geometrica/}{serie geometrica}.

\section{Confronto e confronto asintotico}\label{confronto-e-confronto-asintotico}

Se \(0\le a_n\le b_n\) definitivamente:

\[
\displaystyle
\sum b_n\ \text{converge}
\Longrightarrow
\sum a_n\ \text{converge},
\]

\[
\displaystyle
\sum a_n\ \text{diverge}
\Longrightarrow
\sum b_n\ \text{diverge}.
\]

Se entrambe convergono:

\[
\displaystyle
0\le
R_N^{(a)}
=
\sum_{n=N+1}^{\infty}a_n
\le
\sum_{n=N+1}^{\infty}b_n
=
R_N^{(b)}.
\]

La stima vale per ogni \(N\) successivo alla soglia del confronto.

Se \(a_n\ge0\), \(b_n>0\) definitivamente:

\[
\displaystyle
\frac{a_n}{b_n}\to\ell,
\qquad
0<\ell<+\infty
\quad\Longrightarrow\quad
\sum a_n\ \text{e}\ \sum b_n\ \text{hanno lo stesso carattere}.
\]

Casi unilaterali:

\[
\displaystyle
\frac{a_n}{b_n}\to0,\quad
\sum b_n\ \text{converge}
\quad\Longrightarrow\quad
\sum a_n\ \text{converge},
\]

\[
\displaystyle
\frac{a_n}{b_n}\to+\infty,\quad
\sum b_n\ \text{diverge}
\quad\Longrightarrow\quad
\sum a_n\ \text{diverge}.
\]

\section{Criteri del rapporto, della radice e di Raabe}\label{criteri-del-rapporto-della-radice-e-di-raabe}

Se \(a_n\ne0\) definitivamente ed esiste, finito o esteso,

\[
\displaystyle
L=\lim_{n\to\infty}\left\lvert\frac{a_{n+1}}{a_n}\right\rvert,
\]

allora

\[
\displaystyle
\begin{cases}
L<1 &\Longrightarrow \text{convergenza assoluta},\\
L>1 &\Longrightarrow \text{divergenza},\\
L=1 &\Longrightarrow \text{criterio inconcludente}.
\end{cases}
\]

Criterio della radice, valido senza richiedere l'esistenza di un limite ordinario:

\[
\displaystyle
\rho=\limsup_{n\to\infty}\sqrt[n]{\lvert a_n\rvert},
\]

\[
\displaystyle
\begin{cases}
\rho<1 &\Longrightarrow \text{convergenza assoluta},\\
\rho>1 &\Longrightarrow \text{divergenza},\\
\rho=1 &\Longrightarrow \text{criterio inconcludente}.
\end{cases}
\]

Criterio di Raabe, per \(a_n>0\) definitivamente:

\[
\displaystyle
L=\lim_{n\to\infty}
n\left(\frac{a_n}{a_{n+1}}-1\right),
\]

\[
\displaystyle
\begin{cases}
L>1 &\Longrightarrow \sum a_n\ \text{converge},\\
L<1 &\Longrightarrow \sum a_n\ \text{diverge},\\
L=1 &\Longrightarrow \text{criterio inconcludente}.
\end{cases}
\]

Teoria generale: \href{https://ingegnerismo.it/matematica/criteri-convergenza-serie/}{criteri di convergenza per serie}.

\section{Criterio integrale e condensazione}\label{criterio-integrale-e-condensazione}

Se \(f:[N_0,+\infty)\to(0,+\infty)\) è continua e decrescente e \(a_n=f(n)\):

\[
\displaystyle
\sum_{n=N_0}^{\infty}a_n
\ \text{e}\
\int_{N_0}^{+\infty}f(x)\,dx
\quad
\text{hanno lo stesso carattere}.
\]

Se la serie converge, per \(N\ge N_0\):

\[
\displaystyle
\int_{N+1}^{+\infty}f(x)\,dx
\le
R_N
=
\sum_{n=N+1}^{\infty}a_n
\le
\int_N^{+\infty}f(x)\,dx.
\]

Condensazione di Cauchy, per \(a_n\ge0\) decrescente:

\[
\displaystyle
\sum_{n=1}^{\infty}a_n
\ \text{e}\
\sum_{k=0}^{\infty}2^k a_{2^k}
\quad
\text{hanno lo stesso carattere}.
\]

Teoria: \href{https://ingegnerismo.it/matematica/criterio-dell-integrale/}{criterio dell'integrale} e \href{https://ingegnerismo.it/matematica/criterio-di-condensazione-di-cauchy/}{condensazione di Cauchy}.

\section{Serie alternate}\label{serie-alternate}

Se \(a_n\ge0\), \(a_{n+1}\le a_n\) definitivamente e \(a_n\to0\):

\[
\displaystyle
\sum_{n=0}^{\infty}(-1)^n a_n
\quad\text{converge}.
\]

Stima del resto, oltre la soglia di monotonia:

\[
\displaystyle
R_N
=
\sum_{n=N+1}^{\infty}(-1)^n a_n,
\qquad
\lvert R_N\rvert\le a_{N+1}.
\]

Il resto ha inoltre il segno del primo termine omesso, salvo che sia nullo. Per ottenere un errore non superiore a \(\varepsilon\) è sufficiente imporre \(a_{N+1}\le\varepsilon\).

Teoria: \href{https://ingegnerismo.it/matematica/criterio-di-leibniz/}{criterio di Leibniz}.

\section{Sommazione per parti, Dirichlet e Abel}\label{sommazione-per-parti-dirichlet-e-abel}

Posto

\[
\displaystyle
A_n=\sum_{k=0}^{n}a_k,
\qquad
A_{-1}=0,
\]

\[
\displaystyle
\sum_{n=p}^{q}a_nb_n
=
A_qb_q-A_{p-1}b_p
+
\sum_{n=p}^{q-1}A_n(b_n-b_{n+1}).
\]

Dirichlet:

\[
\displaystyle
\sup_n\left\lvert\sum_{k=0}^{n}a_k\right\rvert<+\infty,
\qquad
b_n\ \text{monotona},
\qquad
b_n\to0
\]

\[
\displaystyle
\Longrightarrow
\sum a_nb_n\ \text{converge}.
\]

Abel:

\[
\displaystyle
\sum a_n\ \text{converge},
\qquad
b_n\ \text{monotona e limitata}
\]

\[
\displaystyle
\Longrightarrow
\sum a_nb_n\ \text{converge}.
\]

Teoria: \href{https://ingegnerismo.it/matematica/criteri-di-abel-e-dirichlet/}{criteri di Abel e Dirichlet}.

\section{Prodotto di Cauchy}\label{prodotto-di-cauchy}

\[
\displaystyle
c_n=\sum_{k=0}^{n}a_kb_{n-k}.
\]

Se

\[
\displaystyle
\sum a_n=A,
\qquad
\sum b_n=B,
\]

e almeno una delle due serie converge assolutamente:

\[
\displaystyle
\sum_{n=0}^{\infty}c_n=AB.
\]

Se entrambe convergono assolutamente:

\[
\displaystyle
\sum_{n=0}^{\infty}\lvert c_n\rvert<+\infty.
\]

\section{Convergenza assoluta e condizionata}\label{convergenza-assoluta-e-condizionata}

\[
\displaystyle
\sum\lvert a_n\rvert<+\infty
\quad\Longrightarrow\quad
\sum a_n\ \text{converge}.
\]

Per \(q\ge p\):

\[
\displaystyle
\left\lvert\sum_{n=p}^{q}a_n\right\rvert
\le
\sum_{n=p}^{q}\lvert a_n\rvert.
\]

Per serie reali:

\[
\displaystyle
a_n^+=\frac{\lvert a_n\rvert+a_n}{2},
\qquad
a_n^-=\frac{\lvert a_n\rvert-a_n}{2},
\qquad
a_n=a_n^+-a_n^-.
\]

Convergenza condizionata:

\[
\displaystyle
\sum a_n\ \text{converge},
\qquad
\sum\lvert a_n\rvert\ \text{diverge}.
\]

Una serie assolutamente convergente può essere riordinata arbitrariamente senza cambiare la somma. Per una serie reale condizionatamente convergente, un riordinamento può invece cambiare la somma oppure distruggere la convergenza.

Teoria: \href{https://ingegnerismo.it/matematica/convergenza-assoluta/}{convergenza assoluta e condizionata}.
